\FourSubsection{Translated Periodic Boundary Condition}
\label{SEC:BC:Numerical:Periodic:Inputs:Translated}

Translated periodic boundary conditions are useful for problems where the flow
in one or more directions can be treated as infinite.
%
Examples primarily consist of two-dimensional geometries
(\forexample an airfoil)
that are converted to a three-dimensional body with an infinite span in the
third dimension.
%
Another example would be a 3D channel flow with translated periodic boundary
conditions in both the span-wise \emph{and} streamwise directions.
%
The periodicity in the streamwise direction allows for the flow to recirculate
indefinitely until it becomes a fully developed channel flow.
%
However, the streamwise periodicity requires a body force to drive the
streamwise bulk flow in the absence of a pressure differential generated by
traditional inflow\slash outflow boundary conditions.

%%%%%%%%%%%%%%%%%%%%%%%%%%%%%%%%%%%%%%%%%%%%%%%%%%%%%%%%%%%%%%%%%%%%%%%%%%%%%%%%

For each periodic boundary condition, the input vector
%
\begin{InputExample}
  \item[{\BCinpBC[nmlBC::BCInput/Periodic/Translation][i]%&
                 {Periodic\%Translation(1:3)}%
                 [Periodic\%Translation]}]
\end{InputExample}
%
provides the translation vector (in grid
units) that transforms the current boundary surface to align it exactly with its
corresponding periodic pair.
%
For a given periodic pair with translational periodicity, the following relation
must be satisfied:
%
\begin{InputExample}
  \item[{\BCinptext[i]{Periodic\%Translation(:) = -}
         \BCinptext[j]{Periodic\%Translation(:)}}]
\end{InputExample}

%%%%%%%%%%%%%%%%%%%%%%%%%%%%%%%%%%%%%%%%%%%%%%%%%%%%%%%%%%%%%%%%%%%%%%%%%%%%%%%%

For example, say you have an S-duct that is 5.0 inches long and drops 1.0 inch
vertically, and your grid file for this geometry is also provided in inches.
%
If you wanted to make this problem streamwise periodic, your inputs for boundary
condition $i$ would look like
%
\begin{InputExample}
  \item[{\BCinptext[i]{Periodic\%Translation(1:3) = [ 5.0, 1.0, 0.0]}}]
 %\item[{\BCinptext[i]{Periodic\%Translation(1) = 5.0}}]
 %\item[{\BCinptext[i]{Periodic\%Translation(2) = 1.0}}]
 %\item[{\BCinptext[i]{Periodic\%Translation(3) = 0.0}}]
\end{InputExample}
%
and the inputs for boundary condition $j$ would look like
%
\begin{InputExample}
  \item[{\BCinptext[j]{Periodic\%Translation(1:3) = [-5.0, -1.0, 0.0]}}]
 %\item[{\BCinptext[j]{Periodic\%Translation(1) = -5.0}}]
 %\item[{\BCinptext[j]{Periodic\%Translation(2) = -1.0}}]
 %\item[{\BCinptext[j]{Periodic\%Translation(3) = 0.0}}]
\end{InputExample}
%
Of course, an S-duct is generally not a problem that you would want to make
streamwise periodic.
%
It is merely used here to illustrate an example of the
periodic translational vector having components in multiple directions.

%###############################################################################

The elements on the boundary faces must be identical in order to have one-to-one
connectivity between cell faces.
%
When using translated periodic boundary conditions for large-eddy simulations,
one must make sure the extent of the periodic direction is long enough as to not
artificially constrain the turbulent length scale in the periodic direction.

%%%%%%%%%%%%%%%%%%%%%%%%%%%%%%%%%%%%%%%%%%%%%%%%%%%%%%%%%%%%%%%%%%%%%%%%%%%%%%%%

